\documentclass[a4paper,12pt]{article}

\usepackage[T2A]{fontenc}
\usepackage[utf8]{inputenc}
\usepackage[russian]{babel}
\usepackage{amsmath,amssymb,mathtools}
\usepackage{helvet}
\renewcommand{\familydefault}{\sfdefault}
\usepackage{icomma}
\usepackage{geometry}
\usepackage{microtype}
\usepackage{tikz}
\usetikzlibrary{positioning}
\usepackage{tkz-euclide}
\usepackage{pgfplots}
\pgfplotsset{compat=1.18}
\usepackage{tabularx,booktabs}
\usepackage{enumitem}
\usepackage{hyperref}
\usepackage{parskip}
\usepackage{fancyhdr}
\usepackage{setspace}
\usepackage{xcolor}

\geometry{margin=2.2cm}
\onehalfspacing
\setlength{\headheight}{15pt}

\definecolor{accent}{HTML}{008080}
\definecolor{darkaccent}{HTML}{004D4D}
\definecolor{textgray}{HTML}{333333}
\definecolor{softaccent}{HTML}{DCEEEE}

\hypersetup{
    colorlinks=true,
    linkcolor=darkaccent
}

\pagestyle{fancy}
\fancyhf{}
\fancyhead[L]{Механика. ОГЭ}
\fancyhead[R]{Чек-лист}
\fancyfoot[C]{\thepage}

\tkzSetUpLine[line width=1pt,color=accent]
\tkzSetUpPoint[color=darkaccent,fill=darkaccent,size=3]
\tkzSetUpLabel[color=black,font=\small]

\begin{document}

% =========================================================
% ТИТУЛЬНЫЙ ЛИСТ
% =========================================================

\begin{titlepage}
\centering
\vspace*{2cm}

{\Huge\color{darkaccent} Чек-лист\par}
\vspace{0.7cm}

{\LARGE Механическое движение\\
Основные понятия механики\par}

\vspace{1.5cm}

{\large Подготовка к ОГЭ по физике\par}

\vspace{1cm}

{\large Лёвин Артём Александрович\\
эксклюзивно для Володи Хачатуряна\par}

\vfill

{\large \today\par}

\vspace{1.5cm}

\begin{tikzpicture}
\draw[color=accent,line width=1.2pt]
(0,0) .. controls (2,0.8) and (4,-0.8) .. (8,0.1);
\end{tikzpicture}

\end{titlepage}

% =========================================================
% 1. МЕХАНИЧЕСКОЕ ДВИЖЕНИЕ
% =========================================================

\section*{1. Механическое движение}

\textbf{Механическое движение} --- изменение положения тела
в пространстве относительно других тел с течением времени.

Для описания движения необходимо указать:

\begin{itemize}
    \item какое тело движется;
    \item относительно какого тела рассматривается движение;
    \item в какой момент времени определяется положение.
\end{itemize}

\subsection*{Виды движения}

\begin{tabularx}{\linewidth}{@{}lX@{}}
\toprule
Вид движения & Основная идея\\
\midrule
Поступательное &
Все точки тела движутся одинаково, а любая прямая,
проведённая в теле, остаётся параллельной самой себе.\\

Вращательное &
Все точки тела движутся по окружностям вокруг некоторой оси.\\

Колебательное &
Движение периодически повторяется в противоположных направлениях.\\
\bottomrule
\end{tabularx}

\vspace{0.7cm}

\begin{center}
\begin{tikzpicture}[x=1cm,y=1cm,>=stealth]

% Поступательное
\draw[rounded corners=2pt,color=textgray]
(-0.2,-0.9) rectangle (4.2,1.2);

\node[font=\small\bfseries,color=darkaccent]
at (2,1) {Поступательное};

\draw[fill=softaccent,draw=darkaccent]
(0.5,-0.25) rectangle (1.5,0.35);

\draw[->,color=accent,line width=1.1pt]
(1.7,0.05) -- (3.4,0.05);

\draw[->,color=accent]
(0.8,-0.5) -- (2.5,-0.5);

\draw[->,color=accent]
(1.2,-0.5) -- (2.9,-0.5);

\node[font=\scriptsize,align=center]
at (2,-1.5)
{все точки смещаются\\одинаково};

% Вращательное
\draw[rounded corners=2pt,color=textgray]
(4.7,-1) rectangle (9.1,1.2);

\node[font=\small\bfseries,color=darkaccent]
at (6.9,0.92) {Вращательное};

\draw[color=accent,line width=1pt]
(6.9,0) circle (0.55);

\fill[darkaccent]
(7.45,0) circle (1.5pt);

\draw[->,color=accent,line width=1pt]
(7.45,0)
arc[start angle=0,end angle=285,radius=0.55];

\node[font=\scriptsize,align=center]
at (6.9,-1.5)
{точки движутся\\по окружностям};

% Колебательное
\draw[rounded corners=2pt,color=textgray]
(9.6,-0.9) rectangle (14,1.2);

\node[font=\small\bfseries,color=darkaccent]
at (11.8,0.92) {Колебательное};

\draw[
    color=accent,
    line width=1pt,
    domain=0:360,
    samples=100,
    smooth,
    variable=\t
]
plot
({10.35+1.45*\t/360},
 {0.12+0.18*sin(6*\t)});

\draw[<->,color=darkaccent,line width=1pt]
(10.6,-0.35) -- (13,-0.35);

\node[font=\scriptsize,align=center]
at (11.8,-1.5)
{движение\\повторяется};

\end{tikzpicture}
\end{center}

% =========================================================
% 2. ОСНОВНАЯ ЗАДАЧА МЕХАНИКИ
% =========================================================

\section*{2. Основная задача механики}

Основная задача механики:

\[
\text{определить положение тела в любой момент времени.}
\]

Для этого вводится система отсчёта.

\subsection*{Система отсчёта}

Система отсчёта включает:

\begin{enumerate}
    \item тело отсчёта;
    \item систему координат;
    \item часы.
\end{enumerate}

Например, автомобиль движется по дороге.
За тело отсчёта можно принять начало трассы.

Если автомобиль находится справа от начала координат:

\[
x>0,
\]

если слева:

\[
x<0.
\]

\begin{center}
\begin{tikzpicture}[x=0.9cm,y=0.9cm,>=stealth]

\draw[->,line width=1pt,color=textgray]
(-6,0) -- (6.4,0)
node[right] {$x$};

\foreach \x in {-5,-4,...,5}{
    \draw[color=textgray]
    (\x,0.10)--(\x,-0.10);

    \node[below,font=\scriptsize]
    at (\x,-0.12) {$\x$};
}

\fill[darkaccent]
(0,0) circle (2pt);

\node[above,font=\small]
at (0,0.12)
{тело отсчёта};

\draw[->,color=accent,line width=1.3pt]
(0.5,0.65) -- (4,0.65);

\node[above,font=\small,color=darkaccent]
at (2.25,0.65)
{$x>0$};

\draw[->,color=accent,line width=1.3pt]
(-0.5,1.25) -- (-4,1.25);

\node[above,font=\small,color=darkaccent]
at (-2.25,1.25)
{$x<0$};

\end{tikzpicture}

{\small
Направление оси координат определяет знак координаты.
}
\end{center}

% =========================================================
% 3. МАТЕРИАЛЬНАЯ ТОЧКА
% =========================================================

\section*{3. Материальная точка}

\textbf{Материальная точка} --- тело, размерами которого
в условиях задачи можно пренебречь.

Пример:

\begin{itemize}
    \item поезд при движении между городами можно считать
    материальной точкой;
    \item автомобиль при въезде в гараж часто нельзя считать
    материальной точкой.
\end{itemize}

Главный вопрос:

\[
\text{Сравнимы ли размеры тела с расстоянием его движения?}
\]

\begin{center}
\begin{tikzpicture}[x=1cm,y=1cm,>=stealth]

% Левый пример
\draw[color=textgray]
(-0.2,0) -- (5.3,0);

\draw[fill=softaccent,draw=darkaccent]
(0.4,0.15) rectangle (2.5,0.65);

\draw[fill=softaccent,draw=darkaccent]
(2.65,0.15) rectangle (3.9,0.65);

\foreach \x in {0.8,1.9,3.0,3.6}{
    \draw[fill=white,draw=darkaccent]
    (\x,0.05) circle (0.13);
}

\draw[->,color=accent,line width=1.2pt]
(4.15,0.4) -- (5.1,0.4);

\node[font=\small\bfseries,color=darkaccent]
at (2.5,1.15)
{Между городами};

\node[font=\scriptsize,align=center]
at (2.5,-0.55)
{размер поезда мал\\по сравнению с путём};

% Разделитель
\draw[color=textgray!60]
(6,-0.8) -- (6,1.4);

% Автомобиль и гараж
\draw[color=textgray,line width=1pt]
(7,0) -- (11.9,0);

\draw[color=darkaccent,line width=1pt]
(8.15,0)
-- (8.15,1.2)
-- (9.4,2.0)
-- (10.65,1.2)
-- (10.65,0);

\draw[fill=softaccent,draw=darkaccent]
(8.55,0.18) rectangle (10.15,0.65);

\draw[fill=white,draw=darkaccent]
(8.9,0.12) circle (0.13);

\draw[fill=white,draw=darkaccent]
(9.8,0.12) circle (0.13);

\node[font=\small\bfseries,color=darkaccent]
at (9.4,2.45)
{При въезде в гараж};

\node[font=\scriptsize,align=center]
at (9.4,-0.55)
{размер автомобиля\\существенен};

\end{tikzpicture}
\end{center}

% =========================================================
% 4. ТРАЕКТОРИЯ, ПУТЬ И ПЕРЕМЕЩЕНИЕ
% =========================================================

\section*{4. Траектория, путь и перемещение}

\textbf{Траектория} --- линия, вдоль которой движется тело.

\textbf{Путь} --- длина траектории. Обозначается $l$.

\textbf{Перемещение} --- вектор, соединяющий начальное
и конечное положение тела. Обозначается $\vec{s}$.

Важно:

\[
l\neq |\vec{s}|
\]

в общем случае.

\begin{center}
\begin{tikzpicture}[x=1cm,y=1cm,>=stealth]

\fill[darkaccent]
(0,0) circle (2pt);

\fill[darkaccent]
(8,0.5) circle (2pt);

\node[below left,font=\small]
at (0,0)
{$A$ --- старт};

\node[below right,font=\small]
at (8,0.5)
{$B$ --- финиш};

% Кривая траектория
\draw[color=accent,line width=1.4pt]
(0,0)
.. controls (1.2,2.0) and (2.4,-1.5) .. (3.6,1.1)
.. controls (4.8,2.1) and (6.2,-1.0) .. (8,0.5);

% Перемещение
\draw[->,color=darkaccent,line width=1.2pt]
(0,0) -- (8,0.5);

\node[above,font=\small,color=accent]
at (4.0,1.65)
{траектория};

\node[below,font=\small,color=darkaccent]
at (4.0,0.15)
{$\vec{s}$};

\node[font=\scriptsize,align=center]
at (4,-1.0)
{путь $l$ --- длина кривой; \quad
$|\vec{s}|$ --- расстояние от старта до финиша};

\end{tikzpicture}
\end{center}

\subsection*{Особые случаи}

Если тело совершило полный оборот:

\[
l=2\pi R,\qquad |\vec{s}|=0.
\]

Если тело прошло половину окружности:

\[
l=\pi R,
\]

а модуль перемещения равен диаметру:

\[
|\vec{s}|=2R.
\]

\vspace{0.5cm}

\begin{center}

\begin{minipage}[t]{0.46\linewidth}
\centering
\textbf{Полный оборот}

\vspace{0.3cm}

\begin{tikzpicture}[scale=0.85]

\tkzDefPoint(0,0){O}
\tkzDefPoint(2,0){A}

\tkzDrawCircle(O,A)
\tkzDrawPoints(O,A)

\tkzLabelPoints[below](O)
\tkzLabelPoints[right](A)

% Направление движения
\draw[->,color=accent,line width=1pt]
(1.75,0.95)
arc[start angle=28,end angle=325,radius=1.98];

\node[font=\small,align=center]
at (0,-2.55)
{старт и финиш совпали\\
$l=2\pi R$, \quad $|\vec{s}|=0$};

\end{tikzpicture}
\end{minipage}
\hfill
\begin{minipage}[t]{0.46\linewidth}
\centering
\textbf{Половина окружности}

\vspace{0.3cm}

\begin{tikzpicture}[scale=0.85]

\tkzDefPoint(0,0){O}
\tkzDefPoint(2,0){A}
\tkzDefPoint(-2,0){B}

\tkzDrawCircle(O,A)
\tkzDrawPoints(O,A,B)

% Пройденная половина окружности
\tkzDrawArc[
    color=accent,
    line width=1.3pt
](O,A)(B)

% Перемещение
\tkzDrawSegment[
    ->,
    color=darkaccent,
    line width=1.2pt
](A,B)

\tkzLabelPoints[below](O)
\tkzLabelPoints[right](A)
\tkzLabelPoints[left](B)

\node[font=\small,align=center]
at (0,-2.55)
{$l=\pi R$, \quad $|\vec{s}|=2R$};

\end{tikzpicture}
\end{minipage}

\end{center}

\subsection*{Движение по координатной прямой}

Если тело несколько раз меняет направление,
путь складывается из всех пройденных участков.
Модуль перемещения определяется начальной и конечной координатами.

Рассмотрим движение:

\[
x_0=1\text{ м}
\quad\longrightarrow\quad
x_1=5\text{ м}
\quad\longrightarrow\quad
x_2=2\text{ м}.
\]

\begin{center}
\begin{tikzpicture}[x=1.15cm,y=1cm,>=stealth]

\draw[->,line width=1pt,color=textgray]
(0,0) -- (6.5,0)
node[right] {$x$};

\foreach \x in {0,1,...,6}{
    \draw
    (\x,0.09)--(\x,-0.09);

    \node[below,font=\scriptsize]
    at (\x,-0.12)
    {$\x$};
}

% Точки
\fill[darkaccent]
(1,0) circle (2pt);

\fill[accent]
(5,0) circle (2pt);

\fill[darkaccent]
(2,0) circle (2pt);

% Движение вправо
\draw[->,color=accent,line width=1.3pt]
(1,0.65) -- (5,0.65);

% Возвращение
\draw[->,color=accent,line width=1.3pt]
(5,1.15) -- (2,1.15);

% Перемещение
\draw[<->,color=darkaccent,line width=1pt]
(1,-0.75) -- (2,-0.75);

\node[above,font=\scriptsize]
at (1,0)
{старт};

\node[above,font=\scriptsize]
at (5,0)
{поворот};

\node[below,font=\scriptsize]
at (2,-0.05)
{финиш};

\node[font=\small]
at (3.2,1.65)
{$l=(5-1)+(5-2)=7\text{ м}$};

\node[font=\small]
at (3.2,-1.15)
{$|\vec{s}|=|2-1|=1\text{ м}$};

\end{tikzpicture}
\end{center}

% =========================================================
% 5. ГЕОМЕТРИЯ
% =========================================================

\section*{5. Геометрия в задачах механики}

Физические задачи часто требуют применения геометрии.

\subsection*{Теорема Пифагора}

Для прямоугольного треугольника:

\[
c^2=a^2+b^2,
\]

где $c$ --- гипотенуза.

Следовательно:

\[
c=\sqrt{a^2+b^2}.
\]

\begin{center}
\begin{tikzpicture}[scale=1]

\tkzDefPoint(0,0){A}
\tkzDefPoint(4.8,0){B}
\tkzDefPoint(4.8,3.2){C}

\tkzDrawPolygon(A,B,C)
\tkzDrawPoints(A,B,C)

\tkzMarkRightAngle[size=0.35](A,B,C)

\tkzLabelPoints[below](A,B)
\tkzLabelPoints[above](C)

\tkzLabelSegment[below](A,B){$a$}
\tkzLabelSegment[right](B,C){$b$}
\tkzLabelSegment[above left](A,C){$c$}

\node[font=\small,align=left]
at (8.0,1.55)
{$c$ --- гипотенуза\\
$a$, $b$ --- катеты\\[2pt]
$c=\sqrt{a^2+b^2}$};

\end{tikzpicture}
\end{center}

\subsection*{Как теорема Пифагора связана с перемещением}

Если тело прошло два взаимно перпендикулярных участка
длиной $a$ и $b$, модуль перемещения равен:

\[
|\vec{s}|=\sqrt{a^2+b^2}.
\]

Например, при $a=3$ м и $b=4$ м:

\[
|\vec{s}|=\sqrt{3^2+4^2}=5\text{ м}.
\]

\subsection*{Окружность}

Радиус $R$ --- отрезок от центра окружности
до точки на окружности.

Диаметр:

\[
d=2R.
\]

Длина окружности:

\[
L=2\pi R=\pi d.
\]

\begin{center}
\begin{tikzpicture}[scale=1]

\tkzDefPoint(0,0){O}
\tkzDefPoint(2.5,0){A}
\tkzDefPoint(-2.5,0){B}
\tkzDefPoint(0,2.5){C}

\tkzDrawCircle(O,A)
\tkzDrawSegment(B,A)
\tkzDrawSegment(O,C)

\tkzDrawPoints(O,A,B,C)

\tkzLabelPoints[below](O)
\tkzLabelPoints[right](A)
\tkzLabelPoints[left](B)
\tkzLabelPoints[above](C)

\tkzLabelSegment[below](O,A){$R$}
\tkzLabelSegment[above](B,A){$d=2R$}
\tkzLabelSegment[left](O,C){$R$}

\node[font=\small,align=left]
at (5.6,0.5)
{$L=2\pi R$\\[3pt]
$L=\pi d$};

\end{tikzpicture}
\end{center}

% =========================================================
% 6. АЛГОРИТМЫ
% =========================================================

\section*{6. Алгоритмы решения задач}

\subsection*{Алгоритм определения пути и перемещения}

\begin{enumerate}
    \item Определить начальную и конечную точки движения.
    \item Определить траекторию движения.
    \item Найти длину всей пройденной траектории --- путь.
    \item Соединить начало и конец движения и найти модуль перемещения.
    \item При необходимости использовать теорему Пифагора,
    свойства окружности или координаты.
\end{enumerate}

\begin{center}
\begin{tikzpicture}[>=stealth]

\node[
    draw=accent,
    rounded corners=2pt,
    minimum width=3.0cm,
    minimum height=0.85cm,
    align=center
] (a) at (0,0)
{Начало\\и конец};

\node[
    draw=accent,
    rounded corners=2pt,
    minimum width=3.0cm,
    minimum height=0.85cm,
    align=center
] (b) at (4,0)
{Траектория};

\node[
    draw=accent,
    rounded corners=2pt,
    minimum width=3.0cm,
    minimum height=0.85cm,
    align=center
] (c) at (8,0)
{Путь $l$};

\node[
    draw=accent,
    rounded corners=2pt,
    minimum width=3.2cm,
    minimum height=0.85cm,
    align=center
] (d) at (4,-1.8)
{Перемещение\\$\vec{s}$};

\node[
    draw=textgray,
    rounded corners=2pt,
    minimum width=4.7cm,
    minimum height=0.85cm,
    align=center
] (e) at (9,-1.8)
{Пифагор, окружность,\\координаты};

\draw[->,color=darkaccent]
(a)--(b);

\draw[->,color=darkaccent]
(b)--(c);

\draw[->,color=darkaccent]
(b)--(d);

\draw[->,color=darkaccent]
(d)--(e);

\end{tikzpicture}
\end{center}

\subsection*{Типичные ошибки}

\begin{itemize}
    \item путать путь и перемещение;
    \item забывать, что перемещение может быть равно нулю;
    \item считать путь по прямой вместо длины траектории;
    \item забывать извлекать квадратный корень после применения
    теоремы Пифагора;
    \item путать выражения «на сколько» и «во сколько раз».
\end{itemize}

\subsection*{Мини-памятка}

\[
\boxed{
\begin{aligned}
\text{«на сколько?»} &\ \longrightarrow\ \text{вычитание},\\
\text{«во сколько раз?»} &\ \longrightarrow\ \text{деление}.
\end{aligned}
}
\]

% =========================================================
% ТРЕНИРОВОЧНЫЙ БЛОК
% =========================================================

\newpage

\section*{Тренировочный блок}

\begin{enumerate}[label=\arabic*.]

\item
Автомобиль проехал 120 км за 2 часа.
Можно ли считать автомобиль материальной точкой
при изучении этого движения?
Обоснуйте.

\item
Тело переместилось по прямой из точки
с координатой $x_1=5$ м в точку $x_2=18$ м.
Найдите путь и модуль перемещения.

\item
Турист прошёл по прямой 300 м вперёд
и 100 м назад.
Найдите путь и модуль перемещения.

\item
Человек обошёл круглый парк радиусом 20 м.
Найдите путь и модуль перемещения.

\item
Тело движется по квадрату со стороной 5 м:
из вершины $A$ в вершину $B$, затем в $C$.
Найдите путь и модуль перемещения.

\item
Тело прошло половину окружности радиусом 3 м.
Найдите путь и модуль перемещения.

\item
Два тела брошены вертикально вверх.
Первое достигло высоты 15 м,
второе --- 25 м.
На сколько отличаются их пути?

\item
Человек прошёл по окружности диаметром 8 м
четверть окружности.
Найдите длину пути.

\item
Тело переместилось из точки $A$ в точку $B$,
затем вернулось в точку $C$.
Координаты:
$x_A=2$ м,
$x_B=10$ м,
$x_C=6$ м.
Найдите путь и модуль перемещения.

\item
Тело прошло по двум взаимно перпендикулярным
направлениям 6 м и 8 м.
Найдите модуль перемещения.

\end{enumerate}

% =========================================================
% ОТВЕТЫ
% =========================================================

\newpage

\section*{Ответы}

\begin{center}
\renewcommand{\arraystretch}{1.35}

\begin{tabularx}{\linewidth}{@{}cX@{}}
\toprule
№ & Ответ\\
\midrule

1 &
Можно считать материальной точкой,
если размеры автомобиля несущественны
по сравнению с расстоянием движения.\\

2 &
$l=13$ м,
$|\vec{s}|=13$ м.\\

3 &
$l=400$ м,
$|\vec{s}|=200$ м.\\

4 &
$l=40\pi$ м,
$|\vec{s}|=0$.\\

5 &
$l=10$ м,
$|\vec{s}|=5\sqrt{2}$ м.\\

6 &
$l=3\pi$ м,
$|\vec{s}|=6$ м.\\

7 &
Первый путь: $30$ м,
второй путь: $50$ м.
Разность: $20$ м.\\

8 &
$l=2\pi$ м.\\

9 &
$l=12$ м,
$|\vec{s}|=4$ м.\\

10 &
$|\vec{s}|=\sqrt{6^2+8^2}=10$ м.\\

\bottomrule
\end{tabularx}
\end{center}

\end{document}